\documentclass[UTF8]{ctexart}

\usepackage{xcolor}
\usepackage{listings}
\usepackage{titlesec}
\usepackage[normalem]{ulem}

% 代码块样式：浅灰色背景，无边框
\lstset{
  basicstyle=\ttfamily\small,
  backgroundcolor=\color{gray!10},
  frame=none,
  breaklines=true,
  columns=flexible,
  keepspaces=true,
  showstringspaces=false,
  xleftmargin=0pt,
  xrightmargin=0pt,
  aboveskip=1em,
  belowskip=1em
}

% 行内代码命令：浅灰色背景框
\newcommand{\code}[1]{\colorbox{gray!10}{\texttt{#1}}}

% 标题左对齐
\titleformat{\section}{\normalfont\Large\bfseries\raggedright}{}{0em}{}
\titleformat{\subsection}{\normalfont\large\bfseries\raggedright}{}{0em}{}
\titleformat{\subsubsection}{\normalfont\normalsize\bfseries\raggedright}{}{0em}{}

\begin{document}

\section*{0: 注意事项}

\subsection*{1. 关于编译和运行}

\begin{itemize}
  \item 编译：将你写的 \code{.c} 文件编译为 \code{.exe} 文件
  \item 运行：将你编译得到的 \code{.exe} 文件运行
\end{itemize}

如果你修改了代码但只点了运行，\textbf{你的修改是不会被应用到程序上的}，会导致代码有时候已经能够通过，但本地测试还不通过的问题。

推荐同学们选择编译并运行的快捷键：DEVC++ 为 \code{F11}，VSCode 为 \code{F5}。其他编辑器或者编辑器配置非默认设置的情况可以问 ai 或者网上搜。

\subsection*{2. 关于使用 ai}

在经历了去年期末平均分怒掉一道题的惨案后，助教们决定今年严查 ai。
请各位同学不要在 C 赛使用 ai。同时，毕竟是上机测试，也请减少交头接耳的情况。

关于前三四道题不会做可以问助教，或者查课件。往往前几道题都是课件改编，难度不会很大。
E 赛虽然不做监管，但还是希望同学们独立思考。

\subsection*{3. 关于设备}

请各位同学如果使用游戏本的话 \textbf{一定要带充电线}！因为有些教室插座较少，推荐每个宿舍再带个插线板以防万一。

\subsection*{4. 关于编译错误}

编译错误：代码中存在语法错误，导致编译失败。
常见的错误有忘记加分号，if、for、while 等控制结构后多加了分号等。

根据第一次上机来看，部分学生 \code{printf()} 中的内容没有加双引号，或是加的位置不对，如 \code{printf(Even)} \code{printf("\%d,ai")} 导致编译报错或 RE。

部分同学还将题目给出的数据范围写进了变量的声明，导致编译报错。
同学们可以学习在本地编译时查看编译错误信息进行调试。

\subsection*{5. 关于提交}

\textbf{不推荐直接在 oj 的提交框内直接写答案。} 推荐先在本地编辑器上完成，编译运行且样例通过再进行提交。

个别同学不熟悉复制粘贴的快捷键：复制为 \code{Ctrl+C}，粘贴为 \code{Ctrl+V}。
提交后只有得到 AC 才是通过，出现 WA、CE、PE 等情况请在考试简介处查询是哪里出现了问题。

\section*{1: Hello, World!}

\begin{enumerate}
  \item 题解中出现的中文字符等问题。
  \item 由于祖传课件中的下面这行代码，对部分同学产生了误导，在输出前加上了一个 \code{'\textbackslash n'} 换行，导致输出整体错开了一行。

  \begin{lstlisting}
printf("\nhello, world\n")
  \end{lstlisting}

  错误输出：

  \begin{lstlisting}
（空行）
Hello, World!
（空行）
  \end{lstlisting}

  正确输出：

  \begin{lstlisting}
Hello, World!
（空行）
  \end{lstlisting}

  \item 很多同学没有复制题目内容的习惯，自己手打 Hello world 导致大小写不对 / 缺空格 / 缺逗号 / ...
  \item \code{int mian} 导致 CE。
\end{enumerate}

\section*{2: Miyabi 的修行 2}

\begin{enumerate}
  \item w、c 写反导致输出 0。
  \item 同 1.3，缺标点、拼写错误 ...
\end{enumerate}

\section*{3: 奇数还是偶数}

这题有一个小事故，课上并没有讲关于负数取模的一个性质，即负数取模依旧会得到负数。

如 \code{-8 mod 3 = -2}，导致写 \code{if(n\%2==0)...else...} 的同学通过，但是写 \code{if(n\%2==1)...else...} 的同学却得到了 WA。

因此后面删除了负数的数据点。

\section*{4: Dlyrotz}

\sout{助教玩音游导致的}

主要在于循环的基础运用，使用 while 循环的实现如下：

\begin{lstlisting}
#include<stdio.h>
int main() {
    int i = 1, n;
    scanf("%d", &n);
    printf("<");
    while(i<=n){
        printf("=");
        i = i + 1;
    }
    printf(">");
    return 0;
}
\end{lstlisting}

\section*{5: Glaciaxion}

方法很多，你可以先算出来首项每次加一个公差，也可以选择每次用等差数列通项公式计算第 i 项。

出现了大量的 PE，源于 \code{printf("\%d",...);} 中的 \code{\%d} 后面没有加空格，导致输出的数字挤在了一起。
正确的版本为 \code{printf("\%d ",...);}。

\section*{6: 空城计}

这道题时间限制比较宽松，从 1 开始枚举合法的最小公倍数也可以通过，但还是鼓励大家学习题解的做法。

样例代码：

\begin{lstlisting}
#include <stdio.h>
int main()
{
    int a, b;
    scanf("%d %d", &a, &b);
    if (a == 0 || b == 0)
    {
        printf("0\n");
        return 0;
    }
    int c = 1;
    while (c % a != 0 || c % b != 0)//枚举最小公倍数
    {
        c++;
    }
    printf("%d\n", c);
    return 0;
}
\end{lstlisting}

注意逻辑运算的符号，\code{||} \code{\&\&} \code{!} 分别表示自然语言中的或、且和否定。特别的，\code{!=} 是不等于，单独写一个 \code{0} 为假，\code{1} 及其他任何数字为真。

如：

\begin{lstlisting}
//不会自然终止的循环
while (1)
{
    printf("endless...\n");
}
\end{lstlisting}

\begin{lstlisting}
//寻找第一个a的倍数
int c=a+1;
while (!(c % a != 0))//与c % a == 0等价
{
    c++;
}
\end{lstlisting}

\section*{7: Credits}

\code{\^} 不是幂次运算！

\code{\^} 是位运算中的异或，我们之后会讲到。关于这道题，记录一个 \code{sign = 1}，每次循环 \code{sign = -sign} 即可。具体如何进行幂次运算也会在之后讲到。

\section*{8: muzermat 的挑战赛之旅}

唯一的一个小坑就是初始值为最大值的情况，题解里也有提到。

\section*{9: 多边形判定}

数学题，可以画图试一试。对于编程题，有时候不需要局限于敲代码，不如先从题目推导一下所需的条件。

\section*{10: 是梦}

\sout{提示题} 还是数学题，暴力计算肯定是会超时的。一般的 oj 每秒能够进行 \(10^7\) 到 \(10^8\) 次运算，具体取决于你代码的复杂程度，浮点运算耗时还会更大。

这种题目如果毫无思路的话，还可以试试手动计算前几项猜测规律 \sout{然后蒙一版提交}。

\end{document}